% Consolidated 02-mechanics-dynamics from 2bcf7e8f8b8ab792312bb68085af56aedb51e18b; original TeX fragments preserved.

% === BEGIN 02-mechanics-dynamics-059-070.tex blob 6e78cb895964814152edd3ebadc33e2ce57b235c ===
% Irodov, Arihant PDF pp. 19--21; printed pp. 10--12.
% Section 1.2 Newton dynamics: problems 59--70 with detailed derivations.
\exproblem{59}{10}
一气球及其载荷总质量为 $m$，最初以向下大小为 $w$ 的恒定加速度下降。
忽略空气阻力。求应抛弃多少压舱物，才能使气球改为以同样大小的
加速度 $w$ 向上运动。
\exsolution
设空气浮力大小为 $B$，抛弃前沿向下方向应用牛顿第二定律：
\[
 mg-B=mw\quad\Longrightarrow\quad B=m(g-w).
\]
抛弃质量 $\Delta m$ 后，若浮力在短时间内保持不变，
沿向上方向列式：
\[
 B-(m-\Delta m)g=(m-\Delta m)w.
\]
消去 $B$ 得
\[
 m(g-w)=(m-\Delta m)(g+w)
 \quad\Longrightarrow\quad
 \boxed{\Delta m=\frac{2mw}{g+w}.}
\]
\excheck 要有正浮力且初态正常，需 $0<w<g$。
若浮力随气球体积或高度显著变化，上式仅适用于浮力近似不变的模型。

\exproblem{60}{10}
如原书图1.9，质量 $m_0$ 的重物悬挂于滑轮一端，另一端连接水平面上的
$m_1$，而 $m_1$ 再以另一根轻绳连接水平面上的 $m_2$。
滑轮、绳质量及轴摩擦可忽略；两水平物体与地面的摩擦系数同为 $k$。
题面特别给出 $m_0=m_1=m_2$。求 $m_0$ 向下加速度
以及 $m_1$、$m_2$ 之间细绳的拉力，并讨论可能情形。
\exsolution
先保留三个质量作一般计算，设 $m_0$ 确实向下运动且加速度为 $a$。
系统整体在运动方向列牛顿第二定律：
\[
 m_0g-k(m_1+m_2)g=(m_0+m_1+m_2)a.
\]
故在正在滑动的状态下
\[
 a=\frac{m_0-k(m_1+m_2)}{m_0+m_1+m_2}g.
\]
对 $m_2$ 单独列式，设两水平物体间绳的拉力为 $T_{12}$：
\[
 T_{12}-km_2g=m_2a
 \quad\Longrightarrow\quad
 T_{12}=\frac{(1+k)m_0m_2}{m_0+m_1+m_2}g.
\]
题给三质量均等于 $m$，若 $0\le k<1/2$，得
\[
 \boxed{a=\frac{1-2k}{3}g\quad(m_0\text{ 向下}),\qquad
 T_{12}=\frac{1+k}{3}mg.}
\]
\excheck 若初态静止且 $k\ge1/2$，
静摩擦力足以阻止系统从静止自行起动，故 $a=0$。
在静止状态下，摩擦力可在容许范围内自适应，
$T_{12}$ 未必仅由 $k,m$ 唯一确定；
应区分\emph{运动摩擦模型}和\emph{静止受力平衡}，不能将滑动公式强行沿用。

\exproblem{61}{10--11}
两物块 $1,2$ 在倾角 $\alpha$ 的斜面上首尾相接且共同下滑（图1.10）。
两块质量分别为 $m_1,m_2$，
摩擦系数分别为 $k_1,k_2$，且 $k_1>k_2$；
其中物块1位于下方。求：(a) 滑动时两物块之间的接触力；
(b) 开始沿斜面下滑所需的最小倾角。
\exsolution
取沿斜面向下为正。物块1摩擦力较大，因而上方物块2向下推物块1。
接触力大小记 $N_c$，分别列式：
\[
 \begin{aligned}
 m_1a&=m_1g\sin\alpha-k_1m_1g\cos\alpha+N_c,\\
 m_2a&=m_2g\sin\alpha-k_2m_2g\cos\alpha-N_c.
 \end{aligned}
\]
两式相加求得共同加速度
\[
 a=g\left[\sin\alpha-
 \frac{k_1m_1+k_2m_2}{m_1+m_2}\cos\alpha\right].
\]
代回任一块的方程：
\[
 \boxed{N_c=\frac{m_1m_2}{m_1+m_2}
 (k_1-k_2)g\cos\alpha.}
\]
开始滑动对应共同加速度的临界值零，因此
\[
 \boxed{\tan\alpha_{\min}=
 \frac{k_1m_1+k_2m_2}{m_1+m_2}.}
\]
\excheck $k_1=k_2$ 时 $N_c=0$，
说明它们在相同摩擦模型下不必互相推挤。
临界角公式把系数理解为相应的最大静摩擦系数；
一旦滑动后的力和加速度则按滑动摩擦模型计算。

\exproblem{62}{11}
一小物体沿倾角 $\alpha=15^\circ$ 的斜面向上抛出。
它上升所需时间比滑回出发点的时间短 $\eta=2.0$ 倍，
即 $t_{\downarrow}=\eta t_{\uparrow}$。
求摩擦系数 $k$。
\exsolution
沿斜面上行时，重力和摩擦力均向下，
减速度大小
$a_\uparrow=g(\sin\alpha+k\cos\alpha)$。
设初速为 $v_0$，则
\[
 t_\uparrow=\frac{v_0}{a_\uparrow},\qquad
 s_{\max}=\frac{v_0^2}{2a_\uparrow}.
\]
下行时加速度大小为
$a_\downarrow=g(\sin\alpha-k\cos\alpha)>0$；
从静止返回同一段路程的时间
\[
 t_\downarrow=\sqrt{\frac{2s_{\max}}{a_\downarrow}}
 =\frac{v_0}{\sqrt{a_\uparrow a_\downarrow}}.
\]
所以
\[
 \eta^2=\frac{a_\uparrow}{a_\downarrow}
 =\frac{\sin\alpha+k\cos\alpha}
 {\sin\alpha-k\cos\alpha},
\]
解得
\[
 \boxed{k=\frac{\eta^2-1}{\eta^2+1}\tan\alpha
 =\frac35\tan15^\circ\simeq0.161.}
\]
\excheck 算得 $k<\tan\alpha$，能保证物体从最高点自行滑下；
这一条件不是可以随意忽略的。

\exproblem{63}{11}
斜面角 $\alpha$，质量 $m_1$ 的物体在斜面上、与斜面摩擦系数为 $k$，
轻绳经无摩擦轻滑轮与悬挂物 $m_2$ 相连（图1.11）。
系统初始静止，求何种 $m_2/m_1$ 导致：
(a) $m_2$ 开始下降；(b) $m_2$ 开始上升；(c) 保持静止。
\exsolution
设比值 $r=m_2/m_1$。若 $m_2$ 将下降，
斜面上 $m_1$ 将向上走，摩擦沿斜面向下。
克服静摩擦所需的力学临界条件是
\[
 m_2g>m_1g\sin\alpha+km_1g\cos\alpha.
\]
因此
$\boxed{r>\sin\alpha+k\cos\alpha}$。
若 $m_2$ 将上升，则 $m_1$ 沿斜面下滑，
摩擦反向，条件为
\[
 m_1g\sin\alpha>m_2g+km_1g\cos\alpha,
\]
即 $\boxed{r<\sin\alpha-k\cos\alpha}$。
其他情况下静摩擦力能调整以平衡系统，故
\[
 \boxed{\sin\alpha-k\cos\alpha
 \le r\le \sin\alpha+k\cos\alpha.}
\]
\excheck 若 $\sin\alpha-k\cos\alpha<0$，
则对正质量比 $r>0$，(b) 的上升区域为空。
边界等号代表临界静止，是否启动取决于额外扰动和静摩擦模型。

\exproblem{64}{11}
沿用第63题装置，给定
$\alpha=30^\circ$、$m_2/m_1=2/3$、$k=0.10$，
原先静止的系统释放后，求 $m_2$ 加速度的大小和方向。
\exsolution
下落驱动力相对斜面摩擦阻力的比值为
\[
 r-\sin\alpha-k\cos\alpha
 =\frac23-\frac12-\frac{\sqrt3}{20}
 \simeq0.0801>0.
\]
所以 $m_2$ 向下、$m_1$ 沿斜面向上。
把两个物体合成一个系统：
\[
 m_2g-m_1g\sin\alpha-km_1g\cos\alpha
 =(m_1+m_2)a.
\]
代入得
\[
 \boxed{a=g\,\frac{r-\sin\alpha-k\cos\alpha}{1+r}
 \simeq0.0481g\simeq0.472\,\mathrm{m/s^2},}
\]
方向为 $m_2$ 向下。
\excheck 系统达到滑动阈值，因为
$r>\sin\alpha+k\cos\alpha$，所选摩擦方向自洽。

\exproblem{65}{11}
质量 $m_1$ 的木板放在光滑水平面上，
质量 $m_2$ 的小块放在木板上；对小块施加从零起随时间线性增大的
水平力 $F=at$（此处 $a$ 为正比例系数，而非加速度）。
小块与木板之间摩擦系数为 $k$。求两者加速度
$a_1(t),a_2(t)$ 并作定性图。
\exsolution
在两者尚未相对滑动时，共同加速度为
$A(t)=at/(m_1+m_2)$。
为推动木板所需摩擦力大小是
$f=m_1A(t)=m_1at/(m_1+m_2)$。
当 $f$ 达最大值 $km_2g$ 时开始滑动，所以
\[
 \boxed{t_0=\frac{km_2g(m_1+m_2)}{am_1}.}
\]
在 $0\le t\le t_0$：
\[
 \boxed{a_1=a_2=\frac{at}{m_1+m_2}.}
\]
在 $t\ge t_0$ 的持续相对滑动阶段，
木板受摩擦力 $km_2g$，小块受反向摩擦力，因此
\[
 \boxed{a_1=\frac{km_2g}{m_1},\qquad
 a_2=\frac{at-km_2g}{m_2}.}
\]
加速度曲线在 $t_0$ 处连续：
$a_1$ 此后变为水平线，$a_2$ 此后增长斜率由
$a/(m_1+m_2)$ 增为 $a/m_2$。
\excheck $\lim_{t\downarrow t_0}a_2
=\lim_{t\uparrow t_0}a_1=km_2g/m_1$；
原书图仅给示意，本稿的分段表达式确定了图形。

\exproblem{66}{11}
如图1.12，斜面底边水平长度 $\ell=2.10\,\mathrm m$，
斜面与水平夹角 $\alpha$，滑动摩擦系数 $k=0.140$。
小物体由斜面顶端从静止下滑。求下滑时间最短时的倾角及最短时间。
\exsolution
滑行长度为 $L=\ell/\cos\alpha$，
沿斜面加速度
$a_\parallel=g(\sin\alpha-k\cos\alpha)$。
仅当 $\tan\alpha>k$ 时能自然下滑；所需时间满足
\[
 t^2=\frac{2L}{a_\parallel}
 =\frac{2\ell}{g(\sin\alpha\cos\alpha-k\cos^2\alpha)}.
\]
令分母中的角度函数
$f(\alpha)=\frac12[\sin2\alpha-k(1+\cos2\alpha)]$。
$t$ 最小时 $f$ 最大，求导：
\[
 f'(\alpha)=\cos2\alpha+k\sin2\alpha=0,
\]
在物理允许区间的最大点为
\[
 \boxed{\alpha_{\rm opt}
 =\frac\pi4+\frac12\arctan k
 \simeq48.99^\circ.}
\]
此时 $f_{\max}=(\sqrt{1+k^2}-k)/2$，故
\[
 \boxed{t_{\min}=
 \sqrt{\frac{4\ell}{g(\sqrt{1+k^2}-k)}}
 \simeq0.992\,\mathrm s.}
\]
\excheck 最优角大于 $45^\circ$，因为增大倾角能减小摩擦效应；
极限 $k=0$ 时恢复 $45^\circ$。

\exproblem{67}{11--12}
质量为 $m$ 的物块通过一根绳子沿倾角为 $\alpha$ 的斜面向上拉动
（图1.13）。绳子与斜面夹角为 $\beta$，摩擦系数为 $k$。
求使匀速上拉所需拉力最小时的 $\beta$，以及最小拉力。
\exsolution
沿斜面向上为正。绳拉力 $T$ 减小斜面对物块的正压力，
$N=mg\cos\alpha-T\sin\beta$。
物块匀速上拉时
\[
 T\cos\beta=mg\sin\alpha+kN
 =mg(\sin\alpha+k\cos\alpha)-kT\sin\beta.
\]
故
\[
 T(\beta)=\frac{mg(\sin\alpha+k\cos\alpha)}
 {\cos\beta+k\sin\beta}.
\]
因为
$\cos\beta+k\sin\beta\le\sqrt{1+k^2}$，
等号在 $\tan\beta=k$ 处成立，所以
\[
 \boxed{\beta_{\rm opt}=\arctan k,\qquad
 T_{\min}=\frac{mg(\sin\alpha+k\cos\alpha)}
 {\sqrt{1+k^2}}.}
\]
\excheck 以上假定最优时仍保持斜面接触。
把最优解代入正压力得
$N=mg(\cos\alpha-k\sin\alpha)/(1+k^2)$；
只有 $\cos\alpha\ge k\sin\alpha$ 时才自洽，
否则必须改按接触边界或离面情形分析。

\exproblem{68}{12}
从 $t=0$ 起，质量 $m$ 的小块置于光滑水平面上，
受到大小 $F(t)=at$、与水平方向恒成仰角 $\alpha$ 的拉力
（图1.14），$a>0$。
求物块刚离开水平面那一刻：
(a) 速率；(b) 此前在水平面上移动的路程。
\exsolution
离地以前竖直方向无加速度，支持力
$N(t)=mg-at\sin\alpha$。对于 $0<\alpha<\pi/2$，
刚离地时
\[
 t_*=\frac{mg}{a\sin\alpha}.
\]
水平方向运动满足
$m\ddot x=at\cos\alpha$，
且 $x(0)=\dot x(0)=0$。两次积分：
\[
 v_x(t)=\frac{a\cos\alpha}{2m}t^2,\qquad
 x(t)=\frac{a\cos\alpha}{6m}t^3.
\]
代入 $t_*$：
\[
 \boxed{v_*=\frac{mg^2\cos\alpha}{2a\sin^2\alpha},\qquad
 s_*=\frac{m^2g^3\cos\alpha}{6a^2\sin^3\alpha}.}
\]
\excheck $N(t_*)=0$，此前 $N>0$；
若 $\alpha=0$，竖直分力为零，物块始终不离面，
本式不适用。

\exproblem{69}{12}
质量为 $m$ 的物块在光滑水平面上从静止开始直线运动，
受到大小恒为 $F=mg/3$ 的外力。
外力与水平方向的夹角随已走路程 $s$ 变化为 $\alpha=as$，
其中 $a>0$。求速率 $v$ 关于角度 $\alpha$ 的函数。
\exsolution
仅水平分力做功：
\[
 F_\parallel(s)=\frac{mg}{3}\cos(as).
\]
由动能定理
$\frac12mv^2=\int_0^s F_\parallel(u)\,\mathrm du$，
得到
\[
 \frac12mv^2=\frac{mg}{3a}\sin(as)
 =\frac{mg}{3a}\sin\alpha.
\]
所以在从静止出发并持续向前、$\sin\alpha\ge0$ 的阶段
\[
 \boxed{v(\alpha)=\sqrt{\frac{2g}{3a}\sin\alpha}.}
\]
\excheck 根号不可省略：$g/a$ 的量纲是
$\mathrm{m^2/s^2}$，而速率量纲须为 $\mathrm{m/s}$。
当 $\alpha$ 增至 $\pi/2$ 时速率最大；
之后水平分力向后而做负功。

\exproblem{70}{12}
在水平面上放置一质量为 $2m$ 的物块与一质量为 $m$ 的
电动机—电池组合体，两者之间距离 $\ell$，
地面摩擦系数均为 $k$。绳子一端系于物块，
另一端缠绕在电动机轴上。启动后，物块以恒加速度 $w$
向电动机移动。求两者相遇时间。
\exsolution
设电动机组向物块运动的加速度为 $a_M$，
绳拉力为 $T$。忽略电动机内部质量变化以及绳质量，
两者均朝对方运动，且滑动摩擦力大小分别为 $2kmg$、$kmg$。
对物块：
\[
 T-2kmg=2mw\quad\Longrightarrow\quad T=2m(w+kg).
\]
对电动机组：
\[
 T-kmg=ma_M
 \quad\Longrightarrow\quad
 a_M=2w+kg.
\]
故两者间距以相对加速度 $w+a_M=3w+kg$ 缩短，
初始均静止时有
$\ell=\frac12(3w+kg)t^2$，因此
\[
 \boxed{t=\sqrt{\frac{2\ell}{3w+kg}}.}
\]
\excheck 若书末答案未能清晰显示根号，应以此量纲正确的表达式为准；
$2\ell/(3w+kg)$ 的量纲是时间平方，不能直接等于时间。
% === END 02-mechanics-dynamics-059-070.tex ===

% === BEGIN 02-mechanics-dynamics-071-080.tex blob 3e7e0da71fccca1e57b36c1b6955d551db7289be ===
% Irodov, Problems in General Physics, Arihant PDF pp. 21--23,
% printed pp. 12--14. The figures 1.15--1.23 and the pulley constraints
% were checked against the user's original PDF.
\exproblem{71}{12}
电梯轿厢顶上固定一理想滑轮，轻绳两端分别悬挂质量为 $m_1,m_2$ 的重物。
电梯以恒定加速度 $a_0$ 向上启动。求：（a）$m_1$ 相对电梯和相对井道
的加速度；（b）滑轮施加于轿厢顶部的力。
\exsolution
取竖直向下为电梯相对坐标的正方向。在电梯参考系中等效重力加速度
$g_{\rm eff}=g+a_0$，两端相对加速度大小相等且方向相反：
\[
m_1g_{\rm eff}-T=m_1a'_1,\quad
T-m_2g_{\rm eff}=m_2a'_1 .
\]
于是
\[
\boxed{a'_1=\frac{m_1-m_2}{m_1+m_2}(g+a_0)}
\quad\text{（相对电梯，向下为正）}.
\]
井道中的向下加速度需要减去电梯向上的加速度：
\[
\boxed{a_1=a'_1-a_0
 =\frac{(m_1-m_2)g-2m_2a_0}{m_1+m_2}}.
\]
拉力 $T=2m_1m_2(g+a_0)/(m_1+m_2)$，两股绳均竖直向下拉滑轮，
故滑轮对顶部的向下作用力
\[
\boxed{F=2T=\frac{4m_1m_2}{m_1+m_2}(g+a_0)}.
\]
\excheck 若 $m_1=m_2$，两物体相对电梯静止，井道中的加速度为 $-a_0$
（向下取正），与轿厢共同向上运动。底本书末疑将 $(g+a_0)$ 误印为
$(g-a_0)$；由于题面明确写电梯\emph{向上}加速，本稿按动力学推导并记勘误。

\exproblem{72}{12--13}
图1.15中，物块1位于倾角 $\alpha$ 的光滑斜面上，
物块2悬挂于活动滑轮的另一股细绳上；后者质量是前者的
$\eta$ 倍。滑轮、细绳均理想。求物块2的竖直加速度及可能方向。
\exsolution
令 $m_1=m,\ m_2=\eta m$。图中上方绳跨过定滑轮，
它的另一端悬挂活动滑轮；活动滑轮另绕一根一端固定、
另一端悬挂物块2的轻绳。记物块2向下加速度 $a$。
两次不可伸长的绳长关系表明，活动滑轮相应移动 $a/2$，
物块1沿斜面向上的加速度也为 $a/2$。
若第二股绳张力为 $T$，活动滑轮上端绳张力为 $2T$。故
\[
2T-mg\sin\alpha=\frac12ma,\qquad
\eta mg-T=\eta ma.
\]
消去 $T$，得
\[
\boxed{a=\frac{2g(2\eta-\sin\alpha)}{4\eta+1}}.
\]
\excheck 若 $\eta>\tfrac12\sin\alpha$，物块2向下；
若 $\eta<\tfrac12\sin\alpha$，则 $a<0$，物块2向上；
等号时平衡。公式的符号由绳长关系给出，不能预先固定方向。

\exproblem{73}{13}
如图1.16，$m_0$ 在光滑水平面上，绳经固定滑轮拉着下方活动滑轮；
活动滑轮绕另一根细绳，其两端各挂 $m_1,m_2$。
所有滑轮、细绳质量与摩擦可略。求 $m_1$ 的竖直加速度，并讨论方向。
\exsolution
记下方绳张力为 $T$，则连接 $m_0$ 与活动滑轮的张力为 $2T$。
令 $m_1,m_2$ 的竖直向下加速度分别为 $a_1,a_2$，
$m_0$ 向右加速度为 $a_0$。上绳约束给出活动滑轮向下
加速度 $a_p=a_0$；下绳约束给出 $a_1+a_2=2a_p$。由牛顿第二定律
\[
m_1a_1=m_1g-T,\quad m_2a_2=m_2g-T,\quad
m_0a_0=2T.
\]
代入绳长约束：
\[
2g-T\left(\frac1{m_1}+\frac1{m_2}\right)=\frac{4T}{m_0},
\quad
T=\frac{2gm_0m_1m_2}
 {4m_1m_2+m_0(m_1+m_2)}.
\]
因此
\[
\boxed{a_1=
g\,\frac{4m_1m_2+m_0(m_1-m_2)}
 {4m_1m_2+m_0(m_1+m_2)}}.
\]
\excheck $a_1$ 以向下为正；分子为正、零、负分别对应
$m_1$ 向下、静止、向上。若 $m_1=m_2$，二者均向下拉动活动滑轮，
$m_0$ 被拉向右，不应误当成普通阿特伍德机。

\exproblem{74}{13}
图1.17中长为 $\ell$ 的重杆质量为 $M$，
另一端套在绳上并可滑动的小球质量为 $m<M$，
滑动时绳对小球施加未知恒定摩擦力。
从静止释放后，两物体均作匀加速运动；经过 $t$ 秒，
小球由杆下端所处高度变为与杆上端等高。求摩擦力大小。
\exsolution
设绳对球的摩擦力向上为 $f$。
理想绳中小球处上方的张力即 $f$，故杆的向下加速度
$a_M=g-f/M$，小球的向下加速度 $a_m=g-f/m$。
因杆相对小球向下多移动 $\ell$，有
\[
\ell=\frac12(a_M-a_m)t^2
=\frac{ft^2}{2}\left(\frac1m-\frac1M\right).
\]
所以
\[
\boxed{f=\frac{2\ell mM}{(M-m)t^2}}.
\]
\excheck 右端有力的量纲；题设 $M>m$ 保证 $f>0$。
如果摩擦方向与相对滑动方向不符，给定运动便不能实现。

\exproblem{75}{13}
图1.18的复合滑轮中，小球1的质量为杆2质量的
$\eta=1.8$ 倍，杆长 $\ell=1.00\,\mathrm m$。
小球初始位置与杆下端等高，系统由静止释放。
求小球运动至与杆上端等高所需时间（滑轮、细绳理想）。
\exsolution
以小球1向下为正，利用图示两级滑轮的不可伸长关系，
每一级活动滑轮的位移是由两侧绳段变化之和确定的。
设杆的相对小球加速度为 $a_{\rm rel}$。由各质量的牛顿方程、
活动滑轮上 $T_{\rm upper}=2T_{\rm lower}$ 与绳长约束消元，可得
\[
a_{\rm rel}=\frac{3g(2-\eta)}{4+\eta}.
\]
于是从静止开始，两端的相对高度变化满足
$\ell=\frac12a_{\rm rel}t^2$，所以
\[
\boxed{t=\sqrt{\frac{2\ell(4+\eta)}
 {3g(2-\eta)}}\simeq1.40\,\mathrm s}.
\]
\excheck 只有 $\eta<2$ 才有题述的相遇过程；
$\eta\to2^-$ 时所需时间发散。该结果与本版书末简答相符。

\exproblem{76}{13}
图1.19中物块1与2的质量比为 $\eta=m_1/m_2=4.0$。
物块1起初距地面 $h=0.20\,\mathrm m$，
经活动滑轮与物块2相连。由静止释放，
忽略滑轮、绳质量与摩擦。求物块2所能上升的最大高度。
\exsolution
用物块1向下位移 $x$ 作广义坐标，绳长约束给出
物块2向上位移 $2x$。动能为
$K=\frac12(m_1+4m_2)\dot x^2$，重力做功为
$(m_1-2m_2)gx$，所以碰地前物块1向下加速度为
\[
a_1=\frac{m_1-2m_2}{m_1+4m_2}g
=\frac{\eta-2}{\eta+4}g .
\]
物块1下降 $h$ 后落地，物块2已经上升 $2h$，
其瞬时上升速率满足 $v_2^2=(2\dot x)^2=8a_1h$。
物块1停住后，绳子松弛，物块2继续上升
$\Delta H=v_2^2/(2g)=4a_1h/g$，因此
\[
\boxed{H=2h+\Delta H=\frac{6\eta}{\eta+4}h
=0.60\,\mathrm m}.
\]
\excheck 这是以物块2的\emph{初始高度}为基准的上升量；
物块1触地的非弹性冲击会导致上绳松弛，故不能在触地以后
继续套用共同加速度。

\exproblem{77}{14}
如图1.20，质量为 $m$ 的竖直导杆 $A$ 与质量为
$M=\eta m$ 的光滑楔块 $B$ 相接触。导杆仅能沿竖直方向移动，
楔块只能沿水平面移动，所有接触处均无摩擦。
求两者的加速度大小。
\exsolution
用楔块向右位移 $x$ 作广义坐标。斜面上升倾角为 $\alpha$，
图示接触约束使导杆向下位移为 $y=x\tan\alpha$。
于是
\[
K=\frac12M\dot x^2+\frac12m\dot y^2
=\frac12(M+m\tan^2\alpha)\dot x^2,
\]
而导杆重力做正功 $mg\,\mathrm dy=mg\tan\alpha\,\mathrm dx$。
由拉格朗日方程可得
\[
a_B=\ddot x=\frac{mg\tan\alpha}{M+m\tan^2\alpha}
=\boxed{\frac{g}{\tan\alpha+\eta\cot\alpha}},
\]
\[
a_A=\ddot y=a_B\tan\alpha
=\boxed{\frac{g}{1+\eta\cot^2\alpha}}.
\]
\excheck 楔块向右、导杆向下；令 $\eta\to\infty$，两者
加速度都趋零，符合重楔极限。

\exproblem{78}{14}
图1.21中，质量为 $M$ 的楔块在光滑水平面上滑动，
质量为 $m$ 的物块沿楔块的竖直导向面运动，
其间摩擦系数为 $k$；两只滑轮、细绳均理想，
绳的一端按图固定在地面。求物块相对地面的加速度大小。
\exsolution
由图中绳长不变关系可见，物块的竖直位移大小
等于楔块的水平位移大小：取二者相应加速度分量
为 $(-a,-a)$，故物块相对地面的加速度大小为 $\sqrt2a$。
物块与楔块在水平方向同加速度，其竖直导向面提供
水平支持力 $N=ma$；因物块相对楔块向下滑动，
摩擦阻力竖直向上，大小 $kma$。
对运动系统以楔块水平位移为广义坐标，
动能中的有效质量是 $M+2m$，重力沿相应约束做功为 $mg\,\mathrm dx$，
摩擦耗功为 $kma\,\mathrm dx$，于是
\[
(M+2m)a=mg-kma.
\]
消元得
\[
\boxed{a_{\rm block}=\sqrt2 a
=\frac{\sqrt2\,mg}{M+(2+k)m}}.
\]
\excheck $k=0$ 时结果为 $\sqrt2mg/(M+2m)$；
当 $M\to\infty$ 时物块加速度趋零。斜向约束来自绳路，
不可把图形误读为普通固定楔块斜面下滑。

\exproblem{79}{14}
图1.22中，质量相等的物体1与2分别置于大块 $A$
的水平顶面和竖直侧面，用跨过其固定滑轮的绳相连。
两处摩擦系数均为 $k$，求 $A$ 至少须有多大的水平加速度
才能使两物体相对 $A$ 保持静止。
\exsolution
设 $A$ 向右加速度为 $a$，在随 $A$ 加速的坐标系中，
两物体受向左大小为 $ma$ 的惯性力。
物体1顶面法向力 $N_1=mg$，故可用摩擦
$|f_1|\le kmg$；物体2侧面法向力 $N_2=ma$，
故 $|f_2|\le kma$。
沿顶面和竖直侧面的平衡条件为
\[
T-ma+f_1=0,\qquad T-mg+f_2=0 .
\]
存在满足两处摩擦约束的 $T$ 当且仅当
区间 $[ma-kmg,ma+kmg]$ 与
$[mg-kma,mg+kma]$ 相交。
对 $0\le k<1$，最小可行加速度来自
$ma+kmg=mg-kma$，即
\[
\boxed{a_{\min}=\frac{1-k}{1+k}g}.
\]
\excheck 若 $k\ge1$，即使 $a=0$ 静摩擦也足以
在合适张力下保持平衡，故数学上应取
$a_{\min}=\max\{0,(1-k)g/(1+k)\}$。
方向选错会使两处摩擦上限的作用颠倒。

\exproblem{80}{14}
图1.23中，倾角为 $\alpha$ 的棱柱向左作水平加速，
其上放质量 $m$ 的物块，摩擦系数 $k<\cot\alpha$。
求物块相对棱柱保持静止时所允许的最大加速度 $a$。
\exsolution
在棱柱参考系中，向右的惯性力 $ma$ 在斜面切向
向上分量为 $ma\cos\alpha$，而重力切向向下分量为
$mg\sin\alpha$。斜面法向力为
\[
N=m(g\cos\alpha+a\sin\alpha).
\]
当 $a$ 增大到上滑临界时，静摩擦沿斜面向下且大小为 $kN$：
\[
ma\cos\alpha-mg\sin\alpha=kN.
\]
因此
\[
\boxed{a_{\max}
=g\,\frac{\sin\alpha+k\cos\alpha}
{\cos\alpha-k\sin\alpha}
=g\frac{1+k\cot\alpha}{\cot\alpha-k}}.
\]
\excheck 题设 $k<\cot\alpha$ 保证分母为正。
若 $a$ 较小时物块原本有下滑趋势，
还需校核另一个静摩擦下限；本题只要求上限。
% === END 02-mechanics-dynamics-071-080.tex ===

% === BEGIN 02-mechanics-dynamics-081-090.tex blob 70b5b11d96962f9f0d02cc7da8d34c084909be0c ===
% Irodov, Arihant PDF pp. 23--25, printed pp. 14--16.
\exproblem{81}{14}
图1.23的光滑楔块质量 $m_1$，楔角 $\alpha$；
物块质量 $m_2$，置于斜面上，从静止释放。
楔块置于无摩擦水平面，求楔块水平加速度。
\exsolution
设楔块向右位移为 $X$，物块相对斜面下滑距离为 $s$，
取斜面向右上升，则物块绝对坐标为
$x_2=X-s\cos\alpha,\ y_2=-s\sin\alpha$（忽略常数）。
系统动能
\[
K=\tfrac12(m_1+m_2)\dot X^2
-m_2\dot X\dot s\cos\alpha+\tfrac12m_2\dot s^2,
\]
势能随 $s$ 减少 $m_2gs\sin\alpha$。
水平动量守恒与沿斜面的运动方程为
\[
(m_1+m_2)\ddot X-m_2\ddot s\cos\alpha=0,\quad
\ddot s-\ddot X\cos\alpha=g\sin\alpha.
\]
消去 $\ddot s$ 得
\[
\boxed{a_{\rm wedge}=
\frac{m_2g\sin\alpha\cos\alpha}
{m_1+m_2\sin^2\alpha}},
\]
方向向右（与物块水平运动方向相反）。
\excheck $m_1\to\infty$ 时楔块近似静止；
$\alpha\to0$ 或 $\pi/2$ 时楔块水平加速度趋零。

\exproblem{82}{14}
图1.24中，质量 $M$ 的楔块在无摩擦水平面上运动，
质量 $m$ 的物块在倾角 $\alpha$ 的光滑斜面上滑动。
物块系轻绳，经固定在楔块顶端的轻滑轮，绳的另一端固定于右侧墙壁。
求楔块的加速度。
\exsolution
取楔块向右位移为 $X$。按图的绳长约束，
物块相对楔块沿斜面向下移动的距离也为 $X$；
因此物块对地速度的分量为
\[
v_x=(1-\cos\alpha)\dot X,\qquad
v_y=-\sin\alpha\,\dot X.
\]
物块速率平方为
$2(1-\cos\alpha)\dot X^2$，全系统动能是
\[
K=\frac12\bigl[M+2m(1-\cos\alpha)\bigr]\dot X^2.
\]
物块高度降低 $X\sin\alpha$，重力作功
$mg\sin\alpha\,\mathrm dX$。于是
\[
\boxed{a_M=
\frac{mg\sin\alpha}
{M+2m(1-\cos\alpha)}},
\]
沿水平向右。
\excheck 分母来自物块的对地速度而非单纯的斜面相对速度，
不可误用 $M+m$ 代替。$m\to0$ 时 $a_M\to0$。

\exproblem{83}{14--15}
质量为 $m$ 的质点沿半径 $R$ 的圆周走过四分之一圆周，
求这一时间段内作用力的\emph{平均矢量}的大小：
(a) 匀速率 $v$；(b) 切向加速度恒为 $a_\tau$，
初速度为零。
\exsolution
动量定理给出 $\overline{\bm F}=\Delta\bm p/\Delta t$。
(a) 初末两速度方向相差 $\pi/2$，大小均为 $v$，
故 $|\Delta\bm p|=\sqrt2mv$；
用时 $\Delta t=(\pi R/2)/v$，于是
\[
\boxed{|\overline{\bm F}|=\frac{2\sqrt2\,mv^2}{\pi R}}.
\]
(b) 若初速为零，四分之一圆周的弧长
$\pi R/2=\tfrac12a_\tau(\Delta t)^2$，
末速为 $a_\tau\Delta t$。
所以 $|\Delta\bm p|=ma_\tau\Delta t$，有
\[
\boxed{|\overline{\bm F}|=ma_\tau}.
\]
\excheck 平均\emph{力矢量}的大小不同于
“各时刻力的模再取平均”；后者不可直接用动量定理计算。

\exproblem{84}{15}
质量 $m=70\,\mathrm{kg}$ 的飞行员随飞机沿半径
$R=500\,\mathrm m$ 的竖直圆环以恒速
$v=360\,\mathrm{km/h}$ 飞行。
求飞机环底、环顶及左右侧点处飞行员的视重。
\exsolution
视重指座椅对飞行员作用力的大小。将速度换算为
$v=100\,\mathrm{m/s}$，向心加速度
$v^2/R=20.0\,\mathrm{m/s^2}$。在环底，座椅支持力向上：
\[
N_{\rm bottom}=m\left(g+\frac{v^2}{R}\right)
\simeq2.09\,\mathrm{kN}.
\]
在环顶，指向圆心的方向向下，重力与座椅力同向：
\[
N_{\rm top}=m\left(\frac{v^2}{R}-g\right)
\simeq0.714\,\mathrm{kN}.
\]
在侧点，重力竖直，向心加速度水平，座椅作用力须同时
提供向心分量 $mv^2/R$ 并平衡重力，故
\[
\boxed{N_{\rm side}=
m\sqrt{g^2+(v^2/R)^2}\simeq1.56\,\mathrm{kN}}.
\]
\excheck 顶点要求 $v^2/R\ge g$ 才能维持座椅向心作用，
本题 $20>9.8$，接触条件成立。

\exproblem{85}{15}
长为 $L$ 的轻绳悬挂一小球，初始绳水平、由静止释放。
设摆动过程中绳与向下竖直线夹角为 $\theta$。求：
(a) 加速度大小和绳张力随 $\theta$ 的变化；
(b) 竖直速度分量最大时的张力；
(c) 合加速度恰好水平时的 $\theta$。
\exsolution
能量守恒给出 $v^2=2gL\cos\theta$，
因此法向加速度 $a_n=v^2/L=2g\cos\theta$，
切向加速度的大小 $|a_\tau|=g|\sin\theta|$，于是
\[
\boxed{a=g\sqrt{1+3\cos^2\theta},\quad
T-mg\cos\theta=mv^2/L,\quad T=3mg\cos\theta.}
\]
摆动至最低点过程中，竖直速度大小
$|v_y|=v\sin\theta$，平方正比于
$\cos\theta(1-\cos^2\theta)$，
对 $c=\cos\theta$ 求极值给出 $c=1/\sqrt3$。
故
\[
\boxed{T_{|v_y|\max}=\sqrt3\,mg}.
\]
用径向、切向分解算合加速度的竖直分量：
$a_y=2g\cos^2\theta-g\sin^2\theta
=g(3\cos^2\theta-1)$。
它为零时
\[
\boxed{\theta=\arccos(1/\sqrt3)\simeq54.7^\circ}.
\]
\excheck (b)(c) 的角度恰好相同；两种条件分别由速度极值和
加速度的竖直分量为零给出，不宜只由数值巧合推断。

\exproblem{86}{15}
一单摆在竖直平面内摆动，它在极端位置和最低点的
\emph{加速度大小}相等。求极端位置的摆线与竖直方向的夹角 $\theta_0$。
\exsolution
设摆长 $L$、最高点摆幅 $\theta_0$。
在极端位置速度为零，只有切向重力加速度，
故 $a_{\rm extreme}=g\sin\theta_0$。
最低点切向加速度为零；由能量守恒，
$v_{\rm low}^2=2gL(1-\cos\theta_0)$，故
$a_{\rm low}=2g(1-\cos\theta_0)$。
令两者相等：
\[
\sin\theta_0=2(1-\cos\theta_0)
\Longrightarrow
2\sin(\theta_0/2)\cos(\theta_0/2)
=4\sin^2(\theta_0/2).
\]
舍去平凡解 $\theta_0=0$，得
\[
\boxed{\theta_0=2\arctan\frac12\simeq53.13^\circ}.
\]
\excheck 一般物理摆动要求 $0<\theta_0<\pi$，
此解在允许范围内。

\exproblem{87}{15}
一小物块从光滑球面顶端由静止滑下，球半径为 $R$。
求物块离开球面时，从竖直向上半径量起的角度 $\theta$，
以及离面速率。
\exsolution
能量守恒：$v^2=2gR(1-\cos\theta)$。
沿指向球心的方向列牛顿第二定律
\[
mg\cos\theta-N=\frac{mv^2}{R}.
\]
物块刚离面时 $N=0$，于是
$g\cos\theta=2g(1-\cos\theta)$，得到
\[
\boxed{\cos\theta=\frac23,\quad
\theta\simeq48.19^\circ,\quad
v=\sqrt{\frac{2gR}{3}}}.
\]
\excheck 离面前 $N\ge0$，在 $\cos\theta=2/3$ 恰好变零；
之后不能继续强制沿球面圆周运动。

\exproblem{88}{15}
图1.26中，L 形光滑刚杆水平转动，转轴过 $O$。
质量为 $m$ 的滑套 $A$ 可沿杆的外伸径向段滑动，
由原长为 $\ell_0$、劲度系数 $\kappa$ 的弹簧连到其内侧锚点 $B$。
系统以恒角速度 $\omega$ 绕竖直轴转动。
求弹簧伸长量，说明旋转方向的影响。
\exsolution
记 $OB=b$，则滑套离轴距离
$r=b+\ell_0+\Delta\ell$。
在随杆匀速转动的参考系中，滑套沿径向静止，
弹簧向内的拉力等于离心惯性力：
\[
\kappa\Delta\ell=m\omega^2r
=m\omega^2(b+\ell_0+\Delta\ell).
\]
解得
\[
\boxed{\Delta\ell=
\frac{b+\ell_0}{\kappa/(m\omega^2)-1}},
\qquad \kappa>m\omega^2.
\]
\excheck 弹簧轴向静力平衡只与 $\omega^2$ 有关，
故正转和反转得到相同伸长量。若 $\kappa\le m\omega^2$，
该线性模型不再有有限稳定径向平衡。
图中 $b+\ell_0$ 是弹簧自然长时滑套的离轴半径，
应与弹簧自身原长 $\ell_0$ 区分。

\exproblem{89}{15}
某水平圆形场地半径为 $R$，
地面对车轮的摩擦系数随距圆心距离 $r$ 变化为
$k(r)=k_0(1-r/R)$。
自行车沿同心圆匀速行驶，求能达到最大速率的圆轨迹半径
及其最大速率。
\exsolution
静摩擦必须足以提供向心加速度：
\[
\frac{mv^2}{r}\le k(r)mg
=k_0mg\left(1-\frac rR\right).
\]
因此对任一半径 $r$，
$v^2\le k_0g\,r(1-r/R)$。
右侧为关于 $r$ 的凹二次函数，其顶点在 $r=R/2$，
最大值为 $k_0gR/4$，故
\[
\boxed{r_{\rm opt}=\frac R2,\qquad
v_{\max}=\frac12\sqrt{k_0gR}}.
\]
\excheck 最大速度的公式必须带根号；
不带根号的 $k_0gR/2$ 量纲为速率平方。

\exproblem{90}{16}
汽车沿半径 $R=40\,\mathrm m$ 的水平圆周轨道
以恒定切向加速度 $a_\tau=0.62\,\mathrm{m/s^2}$
从静止启动，车轮与地面的摩擦系数为 $k=0.20$。
求不发生侧滑之前汽车走过的最大路程。
\exsolution
切向加速度恒定，起速为零，
$v^2=2a_\tau s$。
法向加速度为 $a_n=v^2/R=2a_\tau s/R$。
在水平面上，摩擦合力同时承担切向与法向分量，
最大允许加速度为 $kg$，故
\[
a_\tau^2+\left(\frac{2a_\tau s}{R}\right)^2\le (kg)^2.
\]
取等号求首次达到摩擦极限的路程：
\[
\boxed{s_{\max}=\frac R2
\sqrt{\left(\frac{kg}{a_\tau}\right)^2-1}
\simeq60.0\,\mathrm m}.
\]
\excheck 只有 $a_\tau\le kg$ 才能从静止无滑动启动；
本题 $0.62<1.96\,\mathrm{m/s^2}$，条件满足。
% === END 02-mechanics-dynamics-081-090.tex ===

% === BEGIN 02-mechanics-dynamics-091-100.tex blob 50f36fdc18182be8cd133f17eb54fb3578ef3541 ===
% Irodov, Arihant PDF pp. 25--26, printed pp. 16--17.
\exproblem{91}{16}
汽车以恒速沿水平道路行驶，道路中心线为
$y=a\sin(x/\alpha)$，其中 $a,\alpha>0$，
轮胎与地面的摩擦系数为 $k$。
求不发生侧滑时允许的最高匀速车速。
\exsolution
车辆沿平面曲线匀速运动，法向加速度为 $v^2\kappa$，
其中曲率
\[
\kappa(x)=\frac{|y''|}{(1+y'^2)^{3/2}}
=\frac{(a/\alpha^2)|\sin(x/\alpha)|}
{[1+(a/\alpha)^2\cos^2(x/\alpha)]^{3/2}}.
\]
该表达式在 $\cos(x/\alpha)=0$、$|\sin(x/\alpha)|=1$
时取最大值 $\kappa_{\max}=a/\alpha^2$。
摩擦极限 $mv^2\kappa_{\max}\le kmg$ 给出
\[
\boxed{v_{\max}=\alpha\sqrt{\frac{kg}{a}}}.
\]
\excheck $\alpha$ 和 $a$ 都具有长度量纲，所给结果的量纲为速度；
不能遗漏根号。若只走过部分路段，应对该段实际最大曲率求极值。

\exproblem{92}{16}
总质量 $m$、长度 $2\pi R$ 的均匀链条围成半径 $R$ 的圆环，
套在光滑圆锥面上。圆锥半顶角为 $\theta$，
整个链条绕竖直对称轴以角速度 $\omega$ 匀速旋转。
求链条中的张力。
\exsolution
线密度 $\lambda=m/(2\pi R)$。取一极短链段弧长
$\mathrm ds=R\,\mathrm d\varphi$，
两端张力的合力指向圆心，大小
$T\,\mathrm ds/R$。
光滑圆锥的法向支持力大小（每单位长度）为 $n$，
其竖直向上分量为 $n\sin\theta$，
水平径向向外分量为 $n\cos\theta$。
竖直方向平衡：$n\sin\theta=\lambda g$；
径向向心方向：
\[
\frac TR-n\cos\theta=\lambda\omega^2R.
\]
解得
\[
\boxed{T=\frac{mg}{2\pi}
\left(\cot\theta+\frac{\omega^2R}{g}\right)}.
\]
\excheck $\omega=0$ 时链条仍需要张力抵消圆锥法向力的
水平分量，因而张力一般不为零。

\exproblem{93}{16}
轻绳绕过固定滑轮，两端悬挂质量为 $m_1,m_2$ 的重物。
绳与滑轮存在摩擦，质量比 $m_2/m_1=\eta_0>1$ 时刚好开始打滑。
求：（a）绳—轮摩擦系数 $k$；
（b）当质量比增至 $\eta>\eta_0$ 时两重物的加速度。
\exsolution
取与轮接触的微小绳段，其两端张力相差 $\mathrm dT$。
静摩擦临界时切向平衡
$|\mathrm dT|=k\,\mathrm dN$，
而法向平衡给出 $\mathrm dN\simeq T\,\mathrm d\varphi$，
所以 $\mathrm dT/T=k\,\mathrm d\varphi$。
接触半圆包角为 $\pi$，积分得
$T_2/T_1=e^{k\pi}$。临界时两端均静止，
$T_2/T_1=m_2/m_1=\eta_0$，因此
\[
\boxed{k=\frac{\ln\eta_0}{\pi}}.
\]
当 $\eta>\eta_0$，若动摩擦系数沿用同一值，
则滑动期间 $T_2/T_1=\eta_0$。
较重的 $m_2$ 向下、$m_1$ 向上，设加速度 $a$，有
\[
T_1=m_1(g+a),\qquad T_2=m_2(g-a),
\]
结合 $T_2/T_1=\eta_0$，
\[
\boxed{a=g\,\frac{\eta-\eta_0}{\eta+\eta_0}}.
\]
\excheck 若静、动摩擦系数不同，临界比值仅决定静摩擦系数，
滑动加速度不能仅由 $\eta_0$ 唯一确定；本解遵循原书同系数模型。

\exproblem{94}{16}
质量为 $m$ 的质点在半径为 $R$ 的竖直圆筒内壁无摩擦地运动，
初始速率为 $v_0$，速度与水平面成角 $\alpha$。
求质点对筒壁施加的作用力大小。
\exsolution
内壁光滑意味着筒壁法向力始终水平指向筒轴，
重力只改变速度的竖直分量。
因此水平切向速率为常数
$v_\varphi=v_0\cos\alpha$；
圆筒半径固定，所需水平向心加速度为
$v_\varphi^2/R$。由牛顿第二定律
\[
N=\frac{mv_\varphi^2}{R}
=\boxed{\frac{mv_0^2\cos^2\alpha}{R}}.
\]
\excheck 由牛顿第三定律，质点对筒壁的力与筒壁支持力
大小相等、方向相反；它不随竖直落下时的总速率增加而增加。

\exproblem{95}{16}
质量为 $m$ 的质点在 $xy$ 平面内按
$x=a\sin\omega t,\ y=b\cos\omega t$ 运动。
求作用于质点的合力大小及方向。
\exsolution
分别对坐标求两阶导数：
\[
\ddot x=-a\omega^2\sin\omega t=-\omega^2x,\qquad
\ddot y=-b\omega^2\cos\omega t=-\omega^2y.
\]
因此合力矢量
\[
\boxed{\bm F=-m\omega^2(x\bm i+y\bm j)
=-m\omega^2\bm r},
\]
方向总指向坐标原点，大小为
\[
\boxed{|\bm F|=m\omega^2
\sqrt{a^2\sin^2\omega t+b^2\cos^2\omega t}}.
\]
\excheck 轨迹一般为椭圆，但合力并非处处沿椭圆法线；
只是在 $a=b$ 的圆轨迹情形合力总是纯法向。

\exproblem{96}{16}
质量为 $m$ 的物体以初速度矢量 $\bm v_0$ 斜抛，
忽略空气阻力。求：（a）前 $t$ 秒的动量增量；
（b）回到起抛高度的全程动量增量的大小。
\exsolution
动量增量等于重力冲量，在向上取正的直角坐标系中
$\bm g=-g\bm j$：
\[
\boxed{\Delta\bm p(t)=m\bm g\,t=-mgt\,\bm j}.
\]
若初始竖直分量 $v_{0y}>0$，
回到相同高度的飞行时间 $T=2v_{0y}/g$。
故整个飞行期间
\[
\boxed{|\Delta\bm p_{\rm total}|=mgT
=2m v_{0y}=2m v_0\sin\alpha}.
\]
\excheck 水平动量保持不变；此处 $\alpha$ 为初速度
与水平面的正仰角，且默认全程不碰到其他障碍物。

\exproblem{97}{16}
质量为 $m$ 的静止质点在 $0\le t\le\tau$ 期间受时变力
$\bm F(t)=\bm a\,t(\tau-t)$，
其中 $\bm a$ 是恒定矢量，之后外力撤销。
求：（a）外力刚撤销时的动量；
（b）受力期间走过的路程。
\exsolution
由冲量定理
\[
\bm p(t)=\int_0^t\bm a\,u(\tau-u)\,\mathrm du
=\bm a\left(\frac{\tau t^2}{2}-\frac{t^3}{3}\right).
\]
撤力时有
\[
\boxed{\bm p(\tau)=\frac{\bm a\tau^3}{6}}.
\]
由于 $t(\tau-t)\ge0$ 且方向恒定，质点始终沿 $\bm a$
方向运动，故路程等于位移大小：
\[
s=\int_0^\tau\frac{|\bm p(t)|}{m}\,\mathrm dt
=\frac{|\bm a|}{m}\left(\frac{\tau^4}{6}
-\frac{\tau^4}{12}\right)
=\boxed{\frac{|\bm a|\tau^4}{12m}}.
\]
\excheck 参数 $\bm a$ 并非普通加速度，因
$[\bm F]=[\bm a]\,\mathrm s^2$；
动量、路程表达式的量纲均正确。

\exproblem{98}{16}
质量为 $m$ 的质点从静止开始，在固定直线上受
$F(t)=F_0\sin\omega t$ 的外力。求距出发点的
总路程 $s(t)$，说明函数图形形状。
\exsolution
一维运动方程 $m\dot v=F_0\sin\omega t$。
对初速度为零的情形积分得
\[
v(t)=\frac{F_0}{m\omega}(1-\cos\omega t)\ge0.
\]
尽管外力会周期性反向，速度始终非负，
所以路程等于从原点的正向位移：
\[
\boxed{s(t)=\frac{F_0}{m\omega^2}
(\omega t-\sin\omega t)},\quad t\ge0.
\]
\excheck 曲线整体单调不减，每当 $\omega t=2n\pi$
斜率短暂降为零；平均斜率为 $F_0/(m\omega)$。

\exproblem{99}{17}
质量为 $m$ 的质点起初静止，在固定直线上受到
$F(t)=F_0\cos\omega t$ 的作用。
求第一次重新停止的时间、此前路程及最大速率。
\exsolution
积分可得
\[
v(t)=\frac{F_0}{m\omega}\sin\omega t.
\]
首次正时间零点为 $t_*=\pi/\omega$，
在 $0<t<t_*$ 速度非负，因此
\[
\boxed{t_*=\frac\pi\omega,\qquad
s_*=\int_0^{\pi/\omega}v(t)\,\mathrm dt
=\frac{2F_0}{m\omega^2}}.
\]
在 $t=\pi/(2\omega)$ 时 $\sin\omega t=1$，
故
\[
\boxed{v_{\max}=\frac{F_0}{m\omega}}.
\]
\excheck 在 $t_*$ 的瞬间质点停止后会反向运动，
题目所问“首次停止”不包含之后的反向路程。

\exproblem{100}{17}
质量 $m$ 的摩托艇在湖上以速度 $v_0$ 前进，
$t=0$ 时关闭发动机。阻力与速率成正比，
大小为 $rv$，方向与速度相反。
求：（a）停机后运动时间；
（b）速率随已走路程 $s$ 的变化及总路程；
（c）速率降至 $v_0/\eta$（$\eta>1$）期间的平均速率。
\exsolution
取前进方向为正，$m\dot v=-rv$，
解得
\[
\boxed{v(t)=v_0e^{-rt/m}}.
\]
此速度只在 $t\to\infty$ 时趋零，所以在理想模型中
完全停止所需时间为\emph{无限长}。
用 $\dot v=v\,\mathrm dv/\mathrm ds$ 得
$m v\,\mathrm dv/\mathrm ds=-rv$，在 $v>0$ 时可约去 $v$，
积分得到
\[
\boxed{v(s)=v_0-\frac rm s,\qquad
s_{\rm total}=\frac{mv_0}{r}}.
\]
达到 $v_0/\eta$ 的时间是
$t_\eta=(m/r)\ln\eta$，期间位移为
$s_\eta=(mv_0/r)(1-\eta^{-1})$，于是
\[
\boxed{\bar v=\frac{s_\eta}{t_\eta}
=v_0\frac{\eta-1}{\eta\ln\eta}}.
\]
\excheck $v(s)$ 仅在 $0\le s<s_{\rm total}$ 的运动阶段适用；
$0<v(t)\le v_0$，全程路程有限但持续时间无限。
% === END 02-mechanics-dynamics-091-100.tex ===

% === BEGIN 02-mechanics-dynamics-101-110.tex blob 2e5f32be96b05c7ecf9eefe0e322fb96aff2d8af ===
% Irodov, Arihant PDF pp. 26--27, printed pp. 17--18.
\exproblem{101}{17}
子弹穿过厚度为 $h$ 的木板，速率由 $v_0$ 降为 $v$。
若阻力大小与速率平方成正比，求子弹在板内运动的时间。
\exsolution
令阻力为 $F=-kv^2$，则
$m\dot v=-kv^2$。利用
$\dot v=v\,\mathrm dv/\mathrm dx$ 可得
\[
m v\frac{\mathrm dv}{\mathrm dx}=-kv^2
\ \Longrightarrow\
\frac{\mathrm dv}{v}=-\frac km\,\mathrm dx.
\]
从入板到出板积分，
$\ln(v_0/v)=kh/m$。另一方面
\[
\frac{\mathrm d}{\mathrm dt}\left(\frac1v\right)=\frac km,
\quad
t=\frac mk\left(\frac1v-\frac1{v_0}\right).
\]
消去未知 $k/m$：
\[
\boxed{t=\frac{h(v_0-v)}
 {v_0v\ln(v_0/v)}}.
\]
\excheck 当 $v\to v_0$，对数展开使 $t\to h/v_0$，
与几乎无阻力时的通过时间一致。

\exproblem{102}{17}
物块从静止开始沿倾角 $\alpha$ 的斜面下滑，
摩擦系数随已走过的距离 $x$ 增长为 $k(x)=ax$，$a>0$。
求物块直到停止的总路程，以及全过程中的最大速率。
\exsolution
沿斜面方向列方程
$m\dot v=mg(\sin\alpha-ax\cos\alpha)$。
取 $\dot v=v\,\mathrm dv/\mathrm dx$ 并积分初值 $v(0)=0$：
\[
\frac{v(x)^2}{2}
=g\left(x\sin\alpha-\frac12ax^2\cos\alpha\right).
\]
除 $x=0$ 之外，$v=0$ 时的路程为
\[
\boxed{s_{\rm stop}=\frac{2\tan\alpha}{a}}.
\]
速率最大时 $\dot v=0$，即
$x_m=\tan\alpha/a$，代回得
\[
\boxed{v_{\max}=
\sqrt{\frac ga\sin\alpha\tan\alpha}}.
\]
\excheck 在 $x=x_m$ 处加速度由正变负；
在 $s_{\rm stop}$ 前 $v^2\ge0$，能量表达式始终自洽。

\exproblem{103}{17}
质量为 $m$ 的物块静置于水平粗糙面上，摩擦系数为 $k$。
从 $t=0$ 开始施加随时间线性增长的水平力 $F=at$，$a>0$。
求前 $t$ 秒走过的距离。
\exsolution
力尚未超过最大静摩擦力 $kmg$ 时物块静止，
临界时刻
\[
t_0=\frac{kmg}{a}.
\]
当 $t>t_0$，按静、动摩擦系数相同的原书模型，
净水平力为 $a(t-t_0)$，故
\[
m\ddot x=a(t-t_0),\quad
v(t)=\frac{a}{2m}(t-t_0)^2,\quad
s(t)=\frac{a}{6m}(t-t_0)^3.
\]
因此
\[
\boxed{s(t)=
\begin{cases}
0,&0\le t\le t_0,\\[2pt]
\dfrac{a(t-t_0)^3}{6m},&t>t_0.
\end{cases}}
\]
\excheck 在 $t=t_0$ 处，位移、速度、加速度均连续。
若静摩擦系数与动摩擦系数不同，应分别引入两者，
本表达式不再自动适用。

\exproblem{104}{17}
质量 $m$ 的物体由地面以初速 $v_0$ 竖直上抛，
空气阻力总与速度反向，大小为 $kv^2$，$k>0$。
问物体回到原抛出高度时向下的速率 $v'$。
\exsolution
上升阶段取向上位移 $y$、速率 $u$，
重力和阻力均向下：
\[
m u\frac{\mathrm du}{\mathrm dy}=-mg-ku^2.
\]
从 $u=v_0$ 积分到顶点 $u=0$，最高高度为
\[
H=\int_0^{v_0}\frac{mu\,\mathrm du}{mg+ku^2}
=\frac{m}{2k}\ln\left(1+\frac{kv_0^2}{mg}\right).
\]
下落阶段取向下路程 $z$、速率 $u$，
运动方程 $m u\,\mathrm du/\mathrm dz=mg-ku^2$。
从 $u(0)=0$ 得
\[
u(z)^2=\frac{mg}{k}\left(1-e^{-2kz/m}\right).
\]
令 $z=H$，得
\[
\boxed{v'=\frac{v_0}{\sqrt{1+kv_0^2/(mg)}}}.
\]
\excheck 该速率严格小于 $v_0$，且 $k\to0$ 时趋于 $v_0$，
符合空气阻力耗散机械能的事实。

\exproblem{105}{17}
质量为 $m$ 的质点在平面内受大小恒为 $F$ 的力，
力矢量以角速度 $\omega>0$ 匀速旋转。
质点在 $t=0$ 时静止。求：
(a) 速度矢量随时间变化；
(b) 相邻两次停止间的路程及平均速率。
\exsolution
选坐标轴使初始外力沿 $x$ 正向，
$\bm F(t)=F(\cos\omega t\,\bm i+
\sin\omega t\,\bm j)$。由 $\bm v(0)=0$ 积分：
\[
\boxed{\bm v(t)=\frac{F}{m\omega}
[\sin\omega t\,\bm i+(1-\cos\omega t)\bm j]}.
\]
速率为
\[
v(t)=\frac{2F}{m\omega}
\left|\sin\frac{\omega t}{2}\right|.
\]
相邻停止时刻 $t=0$ 与 $t=2\pi/\omega$，
故
\[
\boxed{\Delta s=\int_0^{2\pi/\omega}v(t)\,\mathrm dt
=\frac{8F}{m\omega^2},\qquad
\overline{v}_{\rm speed}=\frac{\Delta s}{2\pi/\omega}
=\frac{4F}{\pi m\omega}}.
\]
\excheck 此处平均\emph{速率}是书末答案中的 $4F/(\pi m\omega)$；
若求平均\emph{速度矢量}，则需位移除以时间，
其值为 $(F/m\omega)\bm j$，二者不相同。

\exproblem{106}{18}
图1.27中，一小圆片在倾角 $\alpha$ 的斜面上滑动，
摩擦系数 $k=\tan\alpha$。
其速度与斜面最大下降方向 $x$ 的夹角为 $\varphi$，
初始 $\varphi_0=\pi/2$，初速为 $v_0$。
求速率 $v$ 关于 $\varphi$ 的表达式。
\exsolution
沿斜面最大下降方向取 $\bm i$，
重力沿斜面分量大小 $g\sin\alpha$，
摩擦加速度大小 $kg\cos\alpha=g\sin\alpha$，
且始终与速度相反。
设 $b=g\sin\alpha$，速度方向单位矢量
$(\cos\varphi,\sin\varphi)$。动力学沿速度切线和法线分解：
\[
\dot v=b(\cos\varphi-1),\qquad
v\dot\varphi=-b\sin\varphi.
\]
消去时间：
\[
\frac{\mathrm dv}{v}
=\frac{1-\cos\varphi}{\sin\varphi}\,\mathrm d\varphi
=\tan\frac\varphi2\,\mathrm d\varphi.
\]
积分并用 $v(\pi/2)=v_0$：
\[
\ln v=-2\ln\cos(\varphi/2)+\ln(v_0/2),
\quad
\boxed{v(\varphi)=\frac{v_0}{1+\cos\varphi}}.
\]
\excheck 初始 $\varphi=\pi/2$ 时恢复 $v_0$；
后续速度逐渐转向下坡方向 $\varphi\to0$，
极限速率为 $v_0/2$。

\exproblem{107}{18}
一条总长 $\ell$、质量均匀的链条铺在光滑球面上，
球半径 $R$，上端固定在顶点，且 $\ell<\pi R/2$。
解除上端约束的瞬间，求链条各元的切向加速度。
\exsolution
设链条从球顶沿同一子午线向下延伸，
取弧长坐标 $s\in[0,\ell]$，球顶对应 $s=0$。
绳不可伸长，各元沿链条的瞬时切向加速度相同，记为 $a$。
沿整条链对切向方向求受力积分，内部拉力彼此抵消，
光滑球面的支持力无切向分量。线密度为
$\lambda=m/\ell$，弧长 $s$ 处重力切向分量为
$\lambda g\sin(s/R)\,\mathrm ds$。于是
\[
ma=\lambda g\int_0^\ell\sin(s/R)\,\mathrm ds
=\lambda gR[1-\cos(\ell/R)].
\]
故
\[
\boxed{a=\frac{gR}{\ell}[1-\cos(\ell/R)]},
\]
所有链元沿子午线向下同加速度。
\excheck $\ell\ll R$ 时
$a\simeq g\ell/(2R)$；链很短时仅有很小的平均切向重力。

\exproblem{108}{18}
一小物块放在半径 $R$ 的光滑球体顶端，
随后球体突然开始以恒定水平加速度 $a_0$ 平移，
物块相对球体滑下。求：
(a) 刚离面时相对球体的速率；
(b) 离面位置半径与竖直向上方向的夹角 $\theta_0$；
并计算 $a_0=g$ 时的角度。
\exsolution
在随球体加速的非惯性系中，物块受重力和恒定水平
惯性力 $m a_0$，方向与球体加速方向相反。
记 $\theta$ 为向\emph{惯性力一侧}的偏角，
坐标为 $x=R\sin\theta$（向惯性力一侧）、$y=R\cos\theta$。
从顶部开始相对静止，保守有效力场的能量式为
\[
\frac{v'^2}{2}=gR(1-\cos\theta)+a_0R\sin\theta.
\]
向球心投影牛顿第二定律：
\[
\frac{v'^2}{R}=g\cos\theta-a_0\sin\theta-\frac Nm.
\]
离面时 $N=0$。联立消去 $v'$：
\[
g\cos\theta_0-a_0\sin\theta_0=\frac23g.
\]
因此
\[
\boxed{v'_{\rm detach}=\sqrt{\frac{2gR}{3}}}
\]
并设 $\eta=a_0/g$，解
$\cos\theta_0-\eta\sin\theta_0=2/3$ 得
\[
\boxed{\cos\theta_0=
\frac{2+\eta\sqrt{5+9\eta^2}}{3(1+\eta^2)}}.
\]
当 $\eta=1$ 时
$\cos\theta_0=(2+\sqrt{14})/6$，
故 $\boxed{\theta_0\simeq16.9^\circ}$。
\excheck $\eta\to0$ 得 $\cos\theta_0=2/3$，
退化为静止球面上的第87题。这里的速率是\emph{相对球}速率。

\exproblem{109}{18}
质点在平面内受一始终垂直于其速度的力，
其大小与到某固定点距离 $r$ 的 $n$ 次方成反比，
$|\bm F|=C/r^n$。求哪些 $n$ 使圆轨道运动稳定。
\exsolution
因 $\bm F\perp\bm v$，速率 $v$ 恒定。
以固定点为极点，圆轨道半径 $R$ 处
$C/R^n=mv^2/R$，即 $C=mv^2R^{n-1}$。
考虑小的径向扰动 $r=R+\delta r$。
垂直速度且指向轨迹凹侧的力，其极坐标分量在线性近似为
\[
F_r\simeq-\frac{C}{r^n},\qquad
F_\varphi\simeq\frac{C}{r^n}\frac{\dot r}{v}.
\]
对角动量应用
$m\,\mathrm d(r^2\dot\varphi)/\mathrm dt=rF_\varphi$，
积分到一阶得到
$\delta(r^2\dot\varphi)=v\,\delta r$，
所以 $\dot\varphi\simeq(v/R)(1-\delta r/R)$。
将这些代入径向运动方程
$m(\ddot r-r\dot\varphi^2)=F_r$，保留一阶项：
\[
\delta\ddot r+\frac{(1-n)v^2}{R^2}\,\delta r=0.
\]
因而稳定的小振动要求恢复系数为正，即
\[
\boxed{n<1}.
\]
\excheck $n=1$ 为线性中性边界，$n>1$ 径向扰动指数增长。
推导已使用“力始终与速度垂直”而非中央力假设，
不能将开普勒中央力稳定性判据直接挪用。

\exproblem{110}{18}
图1.28所示半径为 $R$ 的光滑半圆形竖直圆弧杆
绕过圆心的竖直轴 $OO'$ 以角速度 $\omega$ 恒速旋转。
质量为 $m$ 的小环可在圆弧上自由滑动。
求稳态平衡位置对应的角度 $\theta$，并讨论稳定性。
\exsolution
从圆弧最低点起量角 $\theta$，小环距转轴
$r=R\sin\theta$，高度取 $y=-R\cos\theta$。
在旋转参考系中，沿圆弧的有效势能为
\[
U_{\rm eff}(\theta)=
-mgR\cos\theta-\frac12m\omega^2R^2\sin^2\theta.
\]
平衡位置满足
\[
\frac{\mathrm dU_{\rm eff}}{\mathrm d\theta}
=mR\sin\theta(g-\omega^2R\cos\theta)=0.
\]
故始终有最低点
$\boxed{\theta_1=0}$；
仅当 $\omega^2R\ge g$ 时还有
\[
\boxed{\theta_2=\arccos\frac{g}{\omega^2R}}.
\]
检查有效势能二阶导数：
当 $\omega^2R<g$ 时 $\theta_1$ 稳定；
当 $\omega^2R>g$ 时 $\theta_1$ 失稳而
$\theta_2$ 稳定；临界 $\omega^2R=g$ 处需要考虑高阶项。
\excheck 这是旋转参考系的稳态，而不是小环相对地面静止；
平衡位置随 $\omega^2$ 改变，旋转方向无关。
% === END 02-mechanics-dynamics-101-110.tex ===

% === BEGIN 02-mechanics-dynamics-111-117.tex blob 5bb2483a6a34cc0489c296311eaef76c8b76d368 ===
% Irodov Arihant PDF pp. 27--28, printed pp. 18--19.
% 111--117: Coriolis, rotating frames, Earth's rotation.
\exproblem{111}{18}
在地理纬度 $\varphi=60^\circ$ 处，步枪向正北方的
竖直标线水平射击，子弹初速为 $v=900\,\mathrm{m/s}$，
靶距 $s=1.0\,\mathrm{km}$。
忽略空气阻力，求子弹命中位置偏离标线的水平距离及方向。
\exsolution
取局部坐标 $\bm e_E$ 向东、$\bm e_N$ 向北、$\bm e_U$ 向上，
地球角速度矢量为
$\bm\Omega=\Omega(\cos\varphi\,\bm e_N+
\sin\varphi\,\bm e_U)$。
一阶近似忽略地球自转对弹道时间的影响，
子弹主要速度是 $\bm v=v\,\bm e_N$。
旋转参考系中的科里奥利加速度为
\[
\bm a_{\rm Cor}=-2\bm\Omega\times\bm v
=2\Omega v\sin\varphi\,\bm e_E.
\]
飞行时间 $t\simeq s/v$，故东向偏移
\[
\boxed{x_E\simeq\frac12(2\Omega v\sin\varphi)
\left(\frac sv\right)^2
=\frac{\Omega s^2\sin\varphi}{v}}.
\]
取 $\Omega=7.292\times10^{-5}\,\mathrm{rad/s}$，
得到 $\boxed{x_E\simeq0.070\,\mathrm m=7.0\,\mathrm{cm}}$，
方向\emph{向东}。
\excheck 该题只问相对竖直标线的水平偏差，
不是重力引起的落弹高度。
保留子弹落下造成的次阶东向修正约为
$O(\Omega gt^3)$，远小于主项。

\exproblem{112}{18--19}
水平转盘以恒角速度 $\omega=6.0\,\mathrm{rad/s}$
绕中心竖直轴旋转。
质量 $m=0.50\,\mathrm{kg}$ 的小物体沿转盘直径
以相对转盘恒速 $v'=0.50\,\mathrm{m/s}$ 运动。
当它距转轴 $r=0.30\,\mathrm m$ 时，
求转盘施加于小物体的合接触力。
\exsolution
在随盘旋转坐标系中，物体的相对加速度为零。
水平接触力需补偿离心力和科里奥利力：
\[
\bm F_{\rm horiz}
=-m\omega^2r\,\bm e_r+
2m\omega v'\,\bm e_\varphi ,
\]
其中方位符号随相对运动方向改变，两项恒垂直。
竖直方向接触支持力为 $N_z=mg$。因此盘对物体的
三维接触合力大小为
\[
\boxed{F=m\sqrt{g^2+\omega^4r^2+4\omega^2{v'}^2}}
\simeq7.89\,\mathrm N\approx8.0\,\mathrm N.
\]
\excheck 必须包括竖直法向支持力 $mg$；
只合成离心与科里奥利的水平项会低估题目的“盘施加的合力”。

\exproblem{113}{19}
一根光滑水平杆绕一端 $A$ 以恒定角速度
$\omega=2.00\,\mathrm{rad/s}$ 绕竖直轴转动。
质量 $m=0.50\,\mathrm{kg}$ 的滑套从 $A$ 处开始，
初始沿杆向外的相对速度 $v_0=1.00\,\mathrm{m/s}$。
当它距转轴 $r=0.50\,\mathrm m$ 时，
求随杆参考系中的科里奥利惯性力大小。
\exsolution
滑套的径向运动在旋转坐标系满足
$\ddot r=\omega^2r$；
因为杆光滑，杆没有沿自身的接触分力，
而离心力为 $m\omega^2r$。
把方程乘以 $2\dot r$ 并积分：
\[
\frac{\mathrm d}{\mathrm dt}
(\dot r^2-\omega^2r^2)=0,
\quad
\dot r^2=v_0^2+\omega^2r^2.
\]
科里奥利力垂直于杆，大小
$F_{\rm Cor}=2m\omega|\dot r|$，所以
\[
\boxed{F_{\rm Cor}=
2m\omega\sqrt{v_0^2+\omega^2r^2}
\simeq2.83\,\mathrm N}.
\]
\excheck 杆的实际接触力须提供相应横向加速度；
不能把“杆光滑”误解为其横向支持力为零。

\exproblem{114}{19}
一水平圆盘半径为 $R$，绕盘缘的竖直固定轴
以角速度 $\omega$ 转动。
质量为 $m$ 的粒子沿盘周以相对盘恒定速率运动。
在粒子离转轴最远时，随盘参考系中的惯性力合力恰为零。
求：（a）粒子的相对盘向心加速度；
（b）惯性力合力大小随其离转轴距离 $r$ 的关系。
\exsolution
取转轴为原点，盘心在 $(R,0)$。
盘周粒子相对盘的位置可写为
$\bm r=R(1+\cos\theta,\sin\theta)$，
距转轴 $r=2R|\cos(\theta/2)|$。
在随盘参考系中，
\[
\bm F_{\rm in}=
m\omega^2\bm r-2m\bm\Omega\times\bm v'.
\]
最远点 $\theta=0$、$r=2R$ 时，惯性力要抵消，
必须沿盘周\emph{逆于盘转向}运动，且满足
$2m\omega v'=m\omega^2(2R)$，即
$v'=\omega R$。因此相对盘沿半径 $R$ 圆周的向心加速度为
\[
\boxed{a'_n=\frac{{v'}^2}{R}=\omega^2R}.
\]
将相对盘的速度取为
$\bm v'=-\omega R(-\sin\theta,\cos\theta)$，
代回惯性力表达式，整理为
$\bm F_{\rm in}
=m\omega^2R(1-\cos\theta,-\sin\theta)$，其大小
\[
\boxed{|\bm F_{\rm in}(r)|=
m\omega^2\sqrt{4R^2-r^2}}.
\]
\excheck 当 $r=2R$ 合力为零；
当 $r=0$，其大小为 $2m\omega^2R$，
与盘心移动及科里奥利项的矢量合成一致。

\exproblem{115}{19}
质量为 $m=0.30\,\mathrm{kg}$ 的小物块从光滑球体
顶部由静止开始滑下，球半径 $R=1.00\,\mathrm m$，
球体以恒角速度 $\omega=6.0\,\mathrm{rad/s}$
绕经过球心的竖直轴旋转。
在随球参考系中，求离面瞬间离心惯性力与科里奥利惯性力大小。
\exsolution
因为球面光滑且轴对称，物块在惯性系中不受绕竖直轴的切向力；
从顶点释放时方位角动量为零，故物块沿固定的竖直子午面滑动，
其离面条件与不转动球的第87题相同：
\[
\cos\theta=\frac23,\quad
\sin\theta=\frac{\sqrt5}{3},\quad
v^2=\frac{2gR}{3}.
\]
离轴水平距离为
$r=R\sin\theta=R\sqrt5/3$，所以离心力
\[
\boxed{F_{\rm cf}=m\omega^2r
=\frac{\sqrt5}{3}m\omega^2R\simeq8.05\,\mathrm N}.
\]
物块在旋转系中的相对速度
$\bm v'=\bm v-\bm\Omega\times\bm r$。
其水平分量有径向 $v\cos\theta$ 与方位向
$-\omega r$ 两个互相垂直的分量，
所以
\[
F_{\rm Cor}=2m\omega\sqrt{v^2\cos^2\theta+\omega^2r^2}
=\boxed{\frac{2m\omega R}{3}
\sqrt{5\omega^2+\frac{8g}{3R}}}
\simeq17.2\,\mathrm N .
\]
\excheck 科里奥利力要用\emph{相对旋转球体}的速度，
不能把惯性系中沿子午线的 $v$ 直接代入 $2m\omega v$。

\exproblem{116}{19}
质量为 $m=2000$ 吨的列车位于北纬 $\varphi=60^\circ$。
求：（a）它以 $v=54\,\mathrm{km/h}$ 沿子午线行驶时，
对钢轨的横向作用力大小和方向；
（b）列车沿哪个方向、以多大速率运动，
可使地球参考系中的离心力与科里奥利惯性力之和为零。
\exsolution
（a）以 $m=2.0\times10^6\,\mathrm{kg}$、
$v=15\,\mathrm{m/s}$ 和
$\Omega=7.292\times10^{-5}\,\mathrm{rad/s}$ 计算，
沿子午线速度产生的科里奥利横向加速度大小是
$2\Omega v\sin\varphi$。列车对轨道的横向作用力
\[
\boxed{F_\perp=2m\Omega v\sin\varphi
\simeq3.79\,\mathrm{kN}}.
\]
在北半球，无论列车向北或向南行驶，科里奥利惯性力
相对于行车方向都向右，因此列车压\emph{右侧钢轨}。
（b）地球自转的离心惯性力矢量与赤道半径的投影成正比，
量级为 $m\Omega^2R_\oplus\cos\varphi$，方向背离地轴。
若列车以速率 $u$ 沿纬线向西行驶，科里奥利惯性力
与该离心力矢量反向；令两者大小相等：
\[
2m\Omega u=m\Omega^2R_\oplus\cos\varphi.
\]
故
\[
\boxed{u=\frac12\Omega R_\oplus\cos\varphi
\simeq116\,\mathrm{m/s}\simeq418\,\mathrm{km/h}},
\]
方向\emph{沿纬线由东向西}。
\excheck （b）比较的是完整的离心和科里奥利惯性力矢量，
不是沿轨道运动时的牵引力。用 $R_\oplus\simeq6371\,\mathrm{km}$
估算；真实地球有效重力与形状修正未包含在理想模型中。

\exproblem{117}{19}
在赤道，某物体相对地面初始静止，
从高度 $h=500\,\mathrm m$ 自由下落，忽略空气阻力。
求它落地时相对铅垂线的水平偏移及方向。
\exsolution
采用当地向东、向北、向上坐标。
赤道处 $\bm\Omega=\Omega\bm e_N$。
在忽略自转的零阶运动中，
竖直速度 $v_z=-gt$，所以
科里奥利加速度的东向分量为
\[
a_E=-2(\bm\Omega\times\bm v)_E=2\Omega gt.
\]
从初始水平速度零积分：
$v_E=\Omega gt^2$、$x_E=\Omega gt^3/3$。
自由落体时间 $t_f=\sqrt{2h/g}$，代入
\[
\boxed{x_E\simeq\frac{\Omega g}{3}
\left(\frac{2h}{g}\right)^{3/2}
=\frac23\Omega h\sqrt{\frac{2h}{g}}
\simeq0.245\,\mathrm m},
\]
方向\emph{向东}。
\excheck 只保留地球自转角速度的一阶效应；
若忽略科里奥利力会错误给出零水平偏差。
% === END 02-mechanics-dynamics-111-117.tex ===
